\section{一次近似}\label{節:一次近似}
\sectionref{節:微分定義}で得た式
\[
 f(a+h)=f(a)+f'(a)h+h\varepsilon(h),\qquad \varepsilon(h)\to0
\]
を使えば,\,$h$ が小さいとき
\[
 f(a+h)\approx f(a)+f'(a)h
\]
と近似できる.これを\textbf{一次近似}という.右辺は $a$ における接線の値である.計算しやすい点 $a$ を近くに選び,そこからの変化を傾きで見積もる.

\begin{example}{平方根の近似}
$\sqrt{4.04}$ を一次近似で求め,近似値が実際の値より大きいか小さいかを確かめよ.
\end{example}
$f(x)=\sqrt{x}$ とすると,\,$f(4)=2$,\,$f'(4)=\frac1{2\sqrt4}=\frac14$ である.\,$4.04=4+0.04$ なので,
\[
 \sqrt{4.04}\approx2+\frac14\cdot0.04=2.01.
\]
$f''(x)=-\cfrac1{4x^{\cfrac32}}<0$ だから,グラフは上に凸で,接線による値は実際の値より大きい.数値によっても,\,$2.01^2=4.0401>4.04$ と確かめられる.

誤差の大きさは,有理化によって直接評価できる.
\[
 2.01-\sqrt{4.04}
 =\frac{2.01^2-4.04}{2.01+\sqrt{4.04}}
 =\frac{0.0001}{2.01+\sqrt{4.04}}.
\]
分母は $4$ より大きいから,\,$0<2.01-\sqrt{4.04}<0.000025$ である.一次近似の式だけでは,指定した小数桁までの正しさが自動的に保証されるわけではない.この例では別の式で誤差を評価した.より一般的な誤差の扱いは\chapterref{章:近似と広義積分}で学ぶ.


